\documentclass[10pt,a4paper]{amsart}
\usepackage[margin=19mm]{geometry}
\usepackage{amsmath,amssymb,mathtools}
\usepackage[scr=rsfs,cal=euler]{mathalfa}
\usepackage{xfrac}
\usepackage[unicode,hidelinks,bookmarks=false]{hyperref}
\newcommand{\ie}{i.e.\ }
\newcommand{\cf}{cf.\ }
\newcommand{\resp}{respectively\ }
\newcommand{\Subsection}{\subsection}
\providecommand{\nobreakdash}{\nobreak\hbox{-}}
\setlength{\emergencystretch}{2em}
\multlinegap0pt
\usepackage{amssymb}
\usepackage{mathtools}
\usepackage[shortlabels]{enumitem}
\usepackage{float}
\usepackage[normalem]{ulem}
\usepackage{xcolor}
\usepackage{tikz}
\newtheorem{lemma}{Lemma}
\newtheorem{claim}{Claim}
\renewcommand{\ge}{\geqslant}
\renewcommand{\geq}{\geqslant}
\renewcommand{\le}{\leqslant}
\renewcommand{\leq}{\leqslant}
\title{Bootstrap percolation on a generalized Hamming cube \MakeUppercase{\romannumeral 2}}
\author{Fengxing Zhu}
\date{Source: arXiv:2503.12607v5 (CC BY 4.0)}
\begin{document}
\maketitle

\noindent\emph{Source and licence.} Bootstrap percolation on a generalized Hamming cube \MakeUppercase{\romannumeral 2}. Version arXiv:2503.12607v5 (CC BY 4.0), distributed by arXiv under the Creative Commons Attribution 4.0 International licence, http://creativecommons.org/licenses/by/4.0/, as recorded in the arXiv licence metadata for this exact version and shown on the abstract page https://arxiv.org/abs/2503.12607v5. Full licence notice: \texttt{licenses/arxiv-2503.12607-CC-BY-4.0.txt}.\par\smallskip

\noindent\textit{Editorial scope.} Counted unit: the article's claim claim (source0.tex lines 471--476). The article's complete original proof of it is delivered with it (source0.tex lines 478--776, one \texttt{proof} environment of 299 lines). No mathematics is written, repaired or re-derived: the statement and the proof are the article's own lines, in the article's own notation, and no retained line is substituted or re-flowed.  Retained uncounted as the source-local context the proof needs: lines 223--228; lines 335--343; lines 369--373.  Nothing was omitted from any retained span, so the package carries no omission record.  No retained line contains a citation, so the wrapper carries no bibliography and no bibliography entry.  No retained line contains a figure environment, an image-inclusion command or drawing code, so no image is delivered and the package declares no asset dependency.  Editorial formatting only, in addition: an AMS \texttt{amsart} preamble that restates the article's own declarations and macros used by the retained lines, namely the environments \texttt{lemma}, \texttt{claim} on separate counters, together with the standard packages those lines need.  The retained environments print in this document as Lemma 1, Claim 1 in order of appearance, whereas the article prints the same items with the numbering of its own full document.  The complete native source of the article as distributed in the arXiv e-print for this exact version is bundled unchanged as \texttt{sources/174867/source0.tex}.
\begin{lemma} \label{lemmaA1}
Let $n \in \mathbb{N}$, $0<p<1$, $t>0$ and $S(n) \sim \text{Binomial}(n,p)$. Then,
$$\mathbb{P}(S(n) \geq np+t) \leq \exp\left(-\frac{2t^2}{n}\right)$$
and
$$\mathbb{P}(S(n) \leq np-t) \leq \exp\left(-\frac{2t^2}{n}\right).$$
\end{lemma}

\begin{lemma} \label{lemma: binom approx}
Suppose $L=N+O(n^{k-1})$ and $Y\sim \text{Binomial}(L,p)$. Suppose an integer $q$ satisfies
$$q-Lp=N\Delta_n+o(N\Delta_n).$$
Then $$
\mathbb P(Y=q) \le N^{-1/2}n^{-\alpha+o(1)},
$$
where $\alpha=\frac{2(\frac{\sqrt{kk!}}{2}-\epsilon)^2}{k!}$.

\end{lemma}

\begin{lemma} Efron-Stein Inequality \label{lemma: efron}

Suppose $X_1,X_2,\cdots,X_n, X_1',X_2', \cdots, X_n'$ are independent with $X_i, X_i'$ having the same distribution for all $i$. Let $X=(X_1,X_2,\cdots,X_n)$ and $X^{(i)}=(X_1,X_2,\cdots, X_{i-1},X_i',X_{i+1},\cdots,X_{n})$. Let $f(X)$ be a function of these random variables. Then 
$$\text{Var}(f(X)) \leq \frac{1}{2}\sum_{i=1}^n \mathbb{E}(f(X)-f(X^{(i)}))^2.$$
\end{lemma}

\begin{claim}
$$
\text{Cov}(I_u,I_w) \le
n^{-\lceil d/2\rceil-2\alpha+o(1)} + \mathbb{I}_{{d\le k}} n^{-k/2-\alpha+o(1)}.
$$
\end{claim}

\begin{proof}
Note that the events $I_u$ and $I_v$ may have dependency due to the common neighbors i.e., $\mathcal{N}(u)\cap \mathcal{N}(w)$ and the facts that $u\in \mathcal{N}(w)$ and $w\in \mathcal{N}(u)$ when $d\le k$.

Let us estimate the Size of the common neighborhood. Let $M:=\mathcal{N}(u)\cap \mathcal{N}(w)$. We first show $M=O\left(n^{k-\lceil d/2\rceil}\right)$.

Without loss of generality, we can assume that $u=[0]^n$, and write  $W:=\{w:||w||_{\ell_1}\}=d$. For a common neighbor $x$, write $a=|x\cap W|$ and $b=|x\setminus W|$.

Then we have 
$$d_H(x,u)=a+b\le k $$
and
$$d_H(x,w)=d-a+b\le k.$$
Adding the two inequalities gives
$$d+2b\le 2k.$$
Therefore,
$$b\le k-\left\lceil\frac d2\right\rceil.$$

There are only $O(1)$ choices for $x\cap W$. Consequently,
$$
M = O_k\left( n^{k-\lceil d/2\rceil}\right).
$$
Note that $M=0$ if $d>2k$.

\textbf{Case 1:} $d >k$.

In this case, $u\notin \mathcal{N}(w)$ and $w\notin \mathcal{N}(u)$. Thus the only initial variables that influence both $I_u$ and $I_w$ are those indexed by the common neighborhood $\mathcal C$.

Let
$$ Z:=\sum_{x\in\mathcal C}\xi_x \sim \text{Binomial}(M,p).
$$

Conditional on all variables $(\xi_x)_{x\in\mathcal C}$, the remaining variables determining $I_u$ and $I_w$ are disjoint. Therefore $I_u$ and $I_w$ are conditionally independent.

Moreover, the conditional probability depends on the common variables only through their sum $Z$. Define
$$
h(s):= \mathbb{P}(I_u=1\mid Z=s).
$$

Since $u$ is initially infected with probability $p$,
$$h(s)= p+(1-p) \mathbb P\left( \text{Binomial}(N-M,p)\ge r-s
\right). $$



The same function $h$ applies to $w$. Conditional independence gives
\begin{align*}
\mathbb{E}[I_uI_w\mid Z]
&= \mathbb{E}[I_u\mid Z]\mathbb{E}[I_w\mid Z] \\
&= h(Z)^2.
\end{align*}
Thus by the tower property of expectation, we have 
$$\text{Cov}(I_u,I_w) =\text{Var}(h(Z)).$$
Therefore it remains to estimate $\text{Var}(h(Z))$.

Write
$$Z=\xi_1+\cdots+\xi_M,$$
where the $\xi_i \sim \text{Bernoulli}(p)$.

Let $\xi_i'$ be an independent copy of $\xi_i$, and let

$$Z^{(i)}= Z-\xi_i+\xi_i'.$$

The Efron–Stein inequality in Lemma~\ref{lemma: efron} gives
$$\text{Var}(h(Z)) \le\frac{1}{2} \sum_{i=1}^{M} \mathbb{E}
\bigl(h(Z)-h(Z^{(i)})\bigr)^2 $$

Let $T_i:=Z-\xi_i \sim \text{Binomial}(M-1,p)$.

If $\xi_i=\xi_i'$, then $Z=Z^{(i)}$. If they differ, then the difference in the arguments of $h$ is $1$. Therefore,

$$\mathbb{E}\bigl(h(Z)-h(Z^{(i)})\bigr)^2= 2p(1-p) \mathbb{E} \bigl(h(T_i+1)-h(T_i)\bigr)^2.
$$

Hence $$\text{Var}(h(Z)) \le Mp(1-p) \mathbb{E}\bigl(h(T+1)-h(T)\bigr)^2 $$
where $T\sim \text{Binomial}(M-1,p)$.

Note that 
$h(s+1)-h(s) =(1-p)\mathbb{P}\left(\text{Binomial}(N-M,p)=r-s-1
\right).$ Now we need to bound this quantity. 

Define the event $$\mathcal{T}= \{T: |T-(M-1)p|\le \sqrt{M} \log n
\}.$$

By Chernoff’s bound in Lemma \ref{lemmaA1}, we have 
$$\mathbb P(\mathcal{T}^c) \le e^{-\Omega(\log^2 n)}.$$

For $T\in\mathcal{T}$, let $q=r-T-1$.Then we have 

\begin{align*}
q-(N-M)p
&= r-T-1-(N-M)p\\
&= \frac{N}{2}-Np+O(1)
-\bigl(T-(M-1)p\bigr)\\
&=N\Delta_n +O(\sqrt M\log n).
\end{align*}

Since $\sqrt{M}\log n=o(N\Delta_n)$, Lemma ~\ref{lemma: binom approx} can be applied. 


Therefore,
$$\mathbb P\left( \text{Binomial}(N-M,p)=q \right)
\le
N^{-1/2}n^{-\alpha+o(1)}$$
uniformly on $\mathcal{T}$.

Therefore, 
$$|h(T+1)-h(T)|
\le
N^{-1/2}n^{-\alpha+o(1)}
$$
on $\mathcal T$. On $\mathcal T^c$, we use the trivial bound $|h(T+1)-h(T)|\le1$. Hence

\begin{align*}
\mathbb E
\bigl(h(T+1)-h(T)\bigr)^2
&\le
N^{-1}n^{-2\alpha+o(1)}
+
e^{-\Omega(\log^2 n)}\\
&=
N^{-1}n^{-2\alpha+o(1)}.
\end{align*}

Consequently, 
$$\text{Var}(h(Z)) \le \frac{M}{N}n^{-2\alpha+o(1)}.$$


Since $N=\Theta(n^k)$,
$$\frac {M}{N} =O_k\left(n^{-\lceil d/2\rceil}\right).$$


Consequently, when $d>k$, 
$$
\text{Cov}(I_u,I_w)
\le
n^{-\lceil d/2\rceil-2\alpha+o(1)}.$$

This is the first term in (1).

\textbf{Case 2:} $d\leq k$

Now $u$ and $w$ are adjacent, so there is additional dependence:
$
w\in \mathcal{N}(u), u\in \mathcal{N}(w).
$

We remove this direct dependence by replacing the two cross-variables with independent copies.

Let $\widehat\xi_u,\widehat\xi_w$ be independent Bernoulli($p$) variables, independent of all original variables. Define
\begin{align*}
J_u &= 
\begin{cases}
1 & \text{if 
$\xi_u=1$ or 
$\sum_{x \in \mathcal{N}(u) \setminus w}\xi_x
+\widehat\xi_w
\ge r$} \\
0 & \text{otherwise}
\end{cases}
\end{align*}

and similarly, 

\begin{align*}
J_w &=
\begin{cases}
1 & \text{if 
$\xi_w=1$
or
$\sum_{x\in \mathcal{N}(w)\setminus{u}}\xi_x
+\widehat\xi_u
\ge r$} \\
0 & \text{otherwise}
\end{cases}
\end{align*}

Note that  $J_u\overset d=I_u,
\qquad
J_w\overset d=I_w.$


After this replacement, the only variables shared by $J_u$ and $J_w$ are the common-neighbor variables $(\xi_x)_{x\in\mathcal C}$.

Indeed, put $U=\mathcal{N}(u)\setminus(\mathcal C\cup{w})$ and $W'=\mathcal{N}(w)\setminus(\mathcal C\cup{u})$. Then $U\cap W'=\emptyset$. 

Conditional on the variables in $\mathcal C$, $J_u$ depends on $\xi_u,\widehat\xi_w$, and $(\xi_x)_{x\in U}$ and $J_w$ depends on $\xi_w,\widehat\xi_u$, and $(\xi_x)_{x\in W'}$.

These two collections are disjoint and independent. Therefore, exactly as above,
$$
\text{Cov}(J_u,J_w)
\le
n^{-\lceil d/2\rceil-2\alpha+o(1)}.
$$

It remains to compare $I_u,I_w$ with $J_u,J_w$.

The variables $I_u$ and $J_u$ can differ only when changing the neighbor bit $\xi_w$ to $\widehat\xi_w$ changes whether $u$ reaches threshold.

Let $$ R_u
:=
\sum_{x\in N(u)\setminus{w}}\xi_x
\sim\text{Binomial}(N-1,p).
$$

Then $I_u\ne J_u$ can occur only if all three of the following happen:

$$
\xi_u=0,\qquad
R_u=r-1,\qquad
\xi_w\ne\widehat\xi_w.
$$

In fact, these conditions are also sufficient. Therefore,
\begin{align*}
\mathbb P(I_u\ne J_u)
&=
(1-p)
\mathbb P\left(
\operatorname{Bin}(N-1,p)=r-1
\right)
\mathbb P(\xi_w\ne\widehat\xi_w)\\
&=
2p(1-p)^2
\mathbb P\left(
\operatorname{Binomial}(N-1,p)=r-1
\right).
\end{align*}

Now $r-1-(N-1)p= N\Delta_n+O(1).$ By Lemma ~\ref{lemma: binom approx}, we have
$$
\mathbb P\left(
\operatorname{Binomial}(N-1,p)=r-1
\right)
\le
N^{-1/2}n^{-\alpha+o(1)}.
$$

Since $N^{-1/2}=n^{-k/2+o(1)}$,$$
\mathbb P(I_u\ne J_u)
\le
n^{-k/2-\alpha+o(1)}.
$$


The same bound holds for $w$:
$$
\mathbb P(I_w\ne J_w)
\le
n^{-k/2-\alpha+o(1)}.
$$

Since $J_u=I_u$ and $J_w=I_w$, 
$$
\operatorname{Cov}(I_u,I_w)
-\operatorname{Cov}(J_u,J_w)
=
\mathbb E[I_uI_w]-\mathbb E[J_uJ_w].
$$
If $a,b,c,d \in \{0,1\}$, we have 
$$|ab-cd|
\le |a-c|+|b-d|.$$


Therefore,
\begin{align*}
\left|
\mathbb E[I_uI_w]-\mathbb E[J_uJ_w]
\right|
&\le
\mathbb E|I_u-J_u|
+
\mathbb E|I_w-J_w|\\
&=
\mathbb P(I_u\ne J_u)
+
\mathbb P(I_w\ne J_w)\\
&\le
n^{-k/2-\alpha+o(1)}.
\end{align*}


Therefore, for $d \le k$
$$
\operatorname{Cov}(I_u,I_w)
\le
n^{-\lceil d/2\rceil-2\alpha+o(1)}
+
n^{-k/2-\alpha+o(1)}. 
$$
For $d>k$, the second term is absent. Hence
$$
\operatorname{Cov}(I_u,I_w)
\le
n^{-\lceil d/2\rceil-2\alpha+o(1)}
+
\mathbb{I}_{\{d\le k\}}
n^{-k/2-\alpha+o(1)}. \qedhere
$$

\end{proof}

\end{document}
